% ============================================================================
% Chapitre 5 — Étude des technologies utilisées
% ============================================================================
\chapter{Étude des technologies utilisées}

\section*{Introduction}

Chaque technologie intégrée dans \caimp{} a fait l'objet d'une évaluation
systématique : pertinence fonctionnelle, maturité, communauté, empreinte
mémoire, et cohérence avec l'ensemble du système. Ce chapitre présente et
justifie chacun de ces choix.

\section{Langages de programmation}

\subsection{Python 3.12}

Python est le langage dominant pour les applications d'intelligence artificielle,
notamment grâce à l'écosystème NumPy/SciPy et aux bibliothèques de traitement
de texte. Dans \caimp{}, il est utilisé pour les services \textsc{api} (FastAPI),
les workers asynchrones (AI Worker, Alert Engine, Forecast Engine, Correlation Engine),
et l'agent de supervision.

La version 3.12 apporte des gains de performance significatifs (10--15~\% sur
les benchmarks CPython) et améliore les messages d'erreur, facilitant le débogage.
La syntaxe \texttt{asyncio} (async/await) est utilisée pervasively pour la
programmation asynchrone, indispensable pour les services qui attendent
simultanément des messages NATS et des réponses Ollama.

\textbf{Alternatives évaluées :} Node.js (moins adapté à NumPy), Ruby on Rails
(absence d'écosystème IA), Java (verbosité excessive pour ce type de service).

\subsection{Go 1.22}

Go a été retenu pour le Telemetry Writer, le service le plus critique en termes
de performance. Son modèle de concurrence basé sur les goroutines et les channels
est idéal pour l'ingestion parallèle de métriques depuis de multiples agents.
Un serveur Go démarre en quelques millisecondes et consomme moins de 20~Mo de RAM
au repos, contre plusieurs centaines pour une application JVM équivalente.

\textbf{Justification spécifique :} le Telemetry Writer traite des batches OTLP
en désérialisant des buffers protobuf, en écrivant en base via \texttt{COPY}
PostgreSQL, et en publiant sur NATS --- trois opérations qui se prêtent
naturellement aux goroutines. Le compilateur Go génère un binaire statique,
simplifiant le packaging Docker.

\subsection{TypeScript (React 18)}

Le frontend est écrit en TypeScript, sur-ensemble typé de JavaScript, compilé
par Vite pour la production. Le typage statique est particulièrement précieux
pour un projet comme \caimp{} où les interfaces entre frontend et backend sont
nombreuses et évolutives.

\section{Frameworks backend}

\subsection{FastAPI}

FastAPI \cite{fastapi_docs} est un \textit{framework} web Python asynchrone
basé sur Starlette et Pydantic. Ses avantages décisifs pour \caimp{} :

\begin{itemize}
  \item \textbf{Génération automatique de documentation OpenAPI/Swagger}, accessible
    à \texttt{/docs} pour chaque service ;
  \item \textbf{Validation automatique des payloads} via les modèles Pydantic,
    éliminant le code de validation manuel ;
  \item \textbf{Support natif de async/await}, compatible avec asyncpg et nats-py ;
  \item \textbf{Performances} comparables à Node.js (benchmarks TechEmpower).
\end{itemize}

\textbf{Alternatives évaluées :} Flask (synchrone par défaut, verbeux), Django
REST Framework (trop lourd pour des microservices), Litestar (moins mature).

\subsection{asyncpg}

asyncpg est un pilote PostgreSQL asyncio natif écrit en C, offrant des
performances nettement supérieures à psycopg2 (2 à 3 fois plus rapide en lecture).
Son support du protocole \textit{binary} PostgreSQL élimine la sérialisation
des données en texte, réduisant la latence des requêtes fréquentes.

\section{Base de données}

\subsection{PostgreSQL 16 + TimescaleDB}

PostgreSQL \cite{postgresql16_doc} est le système de gestion de bases de données
relationnelles le plus avancé du monde \textit{open source}. L'extension
TimescaleDB \cite{timescaledb2017} transforme PostgreSQL en base de données
temporelles haute performance par trois mécanismes :

\begin{description}[leftmargin=2cm, style=nextline]
  \item[Hypertables] Partitionnement automatique par tranches de temps
    (chunks) sans modification de l'\textsc{api} SQL. Pour la table \texttt{metrics},
    chaque chunk couvre 7 jours. Les requêtes sur un intervalle de temps ne
    lisent que les chunks pertinents (\textit{chunk exclusion}).
  \item[Compression columnar] Après 7 jours, les chunks sont compressés en format
    columnar (similaire à Parquet). Le taux de compression moyen pour les métriques
    numériques est de 12:1, réduisant 100~Go de données brutes à moins de 9~Go.
  \item[Politiques de rétention] Des politiques automatiques suppriment les
    chunks anciens (par exemple, supprimer les métriques brutes de plus de 90 jours).
\end{description}

\subsection{pgvector}

pgvector \cite{pgvector2021} est une extension PostgreSQL permettant de stocker
et d'interroger des vecteurs d'embedding dans des colonnes de type \texttt{VECTOR(n)}.
\caimp{} l'utilise pour deux types de documents :

\begin{itemize}
  \item Les \textbf{runbooks} : guides de remédiation de l'équipe, indexés avec
    \texttt{nomic-embed-text} (768 dimensions) ;
  \item Les \textbf{rag\_documents} : incidents passés avec leur résolution.
\end{itemize}

L'index \textsc{hnsw} (\textit{Hierarchical Navigable Small World}) est utilisé
pour la recherche approximative de voisins les plus proches par similarité
cosinus. Sa complexité de recherche est $O(\log n)$ contre $O(n)$ pour la
recherche exhaustive, permettant des requêtes sub-milliseconde sur des centaines
de milliers de vecteurs.

\subsection{Redis 7.2}

Redis est utilisé comme cache distribué et pour les fonctionnalités temps réel :

\begin{itemize}
  \item \textbf{Déduplication d'alertes} : clés NX avec TTL de 15 minutes ;
  \item \textbf{Cache API} : résultats de requêtes fréquentes (30 secondes) ;
  \item \textbf{Métriques live} : valeurs brutes des dernières 90 secondes
    pour l'affichage en temps réel (HEC Bridge) ;
  \item \textbf{Liste noire JWT} : jetons de rafraîchissement révoqués.
\end{itemize}

\section{Intelligence artificielle}

\subsection{Ollama}

Ollama est un serveur d'inférence LLM local exposant une \api{} REST compatible
avec l'interface OpenAI. Il supporte les formats GGUF et GGML, permettant
l'exécution de modèles quantifiés sur du matériel standard. \caimp{} l'utilise
pour deux modèles :

\begin{itemize}
  \item \textbf{llama3.1:8b-instruct-q4\_K\_M} : modèle principal pour les
    explications d'anomalies et les narratifs de prévision. La quantification
    Q4\_K\_M (4 bits, méthode K-means) réduit la taille de 16~Go (fp16)
    à 4.7~Go avec une dégradation de qualité inférieure à 2~\% sur les
    benchmarks de raisonnement \cite{dettmers2022quantization} ;
  \item \textbf{nomic-embed-text} : modèle d'embedding de 270~Mo produisant
    des vecteurs de 768 dimensions, optimisé pour la recherche sémantique
    \cite{nomicembed2024}.
\end{itemize}

\subsection{Holt-Winters — Formulation mathématique complète}

Le lissage exponentiel double de Holt-Winters, implémenté en NumPy pur dans
\texttt{forecaster.py}, suit les équations ci-dessous :

\begin{align}
  L_0 &= x_1 \tag{initialisation du niveau} \\
  T_0 &= \frac{x_n - x_1}{n-1} \tag{initialisation de la tendance} \\[6pt]
  L_t &= \alpha \cdot x_t + (1-\alpha)(L_{t-1} + T_{t-1}),
        \quad 0 < \alpha < 1 \tag{mise à jour du niveau} \\[4pt]
  T_t &= \beta(L_t - L_{t-1}) + (1-\beta)T_{t-1},
        \quad 0 < \beta < 1 \tag{mise à jour de la tendance} \\[4pt]
  \hat{x}_{t+h} &= L_t + h \cdot T_t \tag{prévision à h pas}
\end{align}

\noindent Les intervalles de confiance à $\pm 1.28\sigma$ (couvrant environ 80~\%
de la distribution) sont calculés sur les résidus $r_t = x_t - \hat{x}_{t|t-1}$ :
\begin{equation}
  \hat{x}_{t+h} \pm 1.28 \cdot \hat{\sigma}_r
\end{equation}
où $\hat{\sigma}_r = \sqrt{\frac{1}{n}\sum_{t=1}^{n} r_t^2}$ est l'écart-type
empirique des résidus. Le \textit{time to critical} est calculé par interpolation
linéaire sur la séquence des prévisions dès que $\hat{x}_{t+h} \geq 0.9$ (seuil
de 90~\%).

\section{Infrastructure}

\subsection{Docker et Docker Compose v2}

Docker \cite{docker_docs} isole chaque service dans un conteneur immuable,
garantissant la reproductibilité des déploiements. Docker Compose v2 orchestre
les 19 services de \caimp{} avec :

\begin{itemize}
  \item \textbf{Healthchecks} sur chaque service, avec politique de dépendance
    (\texttt{condition: service\_healthy}) pour garantir l'ordre de démarrage ;
  \item \textbf{Profils} Docker pour les composants optionnels
    (\texttt{--profile otel} pour l'OTel Collector, \texttt{--profile full}
    pour Ollama en mode test) ;
  \item \textbf{Volumes nommés} pour la persistance de PostgreSQL, Redis, NATS
    et Ollama ;
  \item \textbf{Réseau bridge} isolé (\texttt{caimp\_net}) : aucun service
    n'est exposé directement sauf via Nginx.
\end{itemize}

\subsection{Nginx}

Nginx sert de reverse proxy centralisé, concentrant les responsabilités
transversales :
\begin{itemize}
  \item Terminaison TLS et redirections HTTPS → HTTP ;
  \item Rate limiting (200~req/min pour les \api{}, 10~req/min pour l'authentification) ;
  \item En-têtes de sécurité (\textsc{csp}, \textsc{hsts}, \texttt{X-Frame-Options}) ;
  \item Proxy SSE sans buffering pour le chat IA (\texttt{proxy\_buffering off}).
\end{itemize}

\subsection{Protocole OTLP}

L'\textsc{otlp} \cite{otlp_spec} est le protocole standard de la CNCF pour
la transmission de métriques, traces et journaux. \caimp{} utilise OTLP/HTTP
avec encoding protobuf entre l'agent et le Telemetry Writer. Le choix de ce
standard ouvert garantit la compatibilité future avec d'autres \textit{backends}
d'observabilité et d'autres langages d'instrumentation.

\section{Sécurité}

\subsection{JWT (JSON Web Token)}

Les jetons JWT \cite{rfc7519} sont utilisés pour l'authentification des requêtes
\api{}. \caimp{} utilise l'algorithme HS256 (HMAC-SHA256) avec une clé secrète
de 64 octets (512 bits), générée par \texttt{openssl rand -hex 64}. Les jetons
d'accès expirent en 15 minutes, limitant la fenêtre d'exploitation en cas
de vol. Les jetons de rafraîchissement (7 jours) sont révocables via la table
\texttt{refresh\_token\_blacklist}.

\subsection{mTLS pour les agents}

L'authentification mutuelle TLS \cite{rfc8446} garantit que seuls les agents
possédant un certificat valide émis par la CA interne peuvent envoyer des métriques.
L'autorité de certification est stockée chiffrée en base (AES-256-GCM) et
n'est jamais exposée via les \api{}. La révocation d'un agent s'effectue en
temps réel via la table \texttt{revoked\_agents}, consultée à chaque connexion.

\section*{Conclusion}

Le choix de chaque technologie dans \caimp{} résulte d'une évaluation rigoureuse
des alternatives. L'association Python/Go pour le backend, TimescaleDB pour les
données temporelles, NATS pour le bus d'événements, et Ollama pour l'inférence
LLM locale constitue un ensemble cohérent, performant et accessible. Le chapitre
suivant détaille l'implémentation concrète de ces choix.
